Table~\ref{tab:generalization-results} separates endpoint generalization from layout and dynamics shifts. The full RL policy changes from 96.0\% success on the training layout with held-out endpoints to 90.0\% on unseen matched-density layouts and 72.0\% on denser layouts. The three classical planners remain at 100\% on these splits, so the experiment provides no reliability advantage for the learned method. Scaling exposes the strongest boundary: RL success is 98.0\%, 56.0\%, 14.0\%, and 0.0\% at grid widths 15, 30, 50, and 100, while every classical planner remains at 100\%. At $100\times100$, the full RL policy uses 932.266 ms of route decision time versus 3.566 ms for A*. Thus, these implementations show no scale crossover in favor of RL.

The held-out-pair ablation in Table~\ref{tab:ablation-results} is also non-monotonic. Vanilla DQN and the full-grid observation reach only 38.0\% and 15.3\%, respectively, supporting Double DQN and the local observation in this setup. Yet removing HER (100.0\%), shaping (100.0\%), or the curriculum (100.0\%) each exceeds the full method's 96.0\% on this split. These one-factor interventions show configuration sensitivity, not separable causal main effects, because the learning components interact.

At matched change events, the full RL method attains 87.9\% post-change success, while its recovery criterion is met in 100.0\% of observed events. This distinction matters: recovery of optimal remaining cost before termination does not guarantee that the complete route succeeds. The corresponding D* Lite values are 100.0\% and 100.0\%.
